\section{An example of time synchronization from National Instruments}

National Instruments is a company that produces test equipment and virtual instrumentation software. One of the most important application for their products is data acquisition, and  they are considered one of the leaders thanks to their experience in this field. 

An article from their website \cite{NISync} states that data should be appropriately synchronized in time to be accurately analyzed, and they propose a list of approaches to fulfill such goal using their products. In particular, they focus on their PXI platform.

\paragraph{The PXI Platform} PXI is defined by National Instruments as the standard for automated test and measurements. It is a family of modular instruments and I/O elements that can work together, with specialized synchronization algorithms and software. The main core of the system is the chassis with its own controller, that accepts a certain number of measurement modules and provides the same functionalities to all of them. 

In the latest version of the platform, called \textit{PXI Express}, developers included different elements thought for precise time synchronization:
\begin{itemize}
    \item A 10~MHz backplane clock, offered in parallel by the chassis to all connected component
    \item A single-ended trigger bus
    \item A star trigger signal, accurately developed such that the connection to all the modules have the same length
    \item 100 MHz differential clock and differential triggers, with increased noise immunity. 
\end{itemize}
The user can then freely choose the way to share clock and trigger signals between all modules. Every option has its own advantages and disadvantages considering signal skew, signal integrity, compatibility with external modules and necessity for external controllers.
For example, the 100~MHz clock provides a skew that is less than 250~ps, while the star trigger signal performs a skew that is less than 500~ps. Both of them provide the best signal integrity between all the solutions. 
Of course, the listed solutions are intended for modules connected to the same chassis, but some components can perform synchronization between different ones by using GPS or IRIG-B standards.  

\paragraph{Synchronization Techniques} In order to effectively perform time synchronization, National Instruments defines the differences between two types of devices:
\begin{itemize}
    \item Sample Clock Timed Devices. These devices use a \textit{sample clock} that controls when the samples are acquired. For example, SAR ADC converters usually have a dedicated convert line that, when toggled, triggers an immediate conversion.
    \item Oversample Clock Timed Devices. In this case, the device requires a free-running oversample clock that drives the conversion while keeping the ADC operating: the resulting samples are generated at a rate that is a division of the original clock. These devices can't be controlled by external signals, hence synchronization is more complicated.
\end{itemize}

To synchronize the sample clock, different options are possible:
\begin{itemize}
    \item Share sample clock directly. A master device will share its sample clock to the slave ones, using one or more of the backplane of the PXI chassis.
    \item Derive the sample clock from a common reference clock. The 100~MHz differential clock is divided to generate the sample clock, distributed to all the devices. A dedicated software from NI is able to perform all the necessary configuration for the user.
    \item Share the trigger signal. This is the simplest but the loosest synchronization mechanism, as all the devices will only receive the same trigger signal. In this case, usually devices drift in time, as sample clocks are not locked together. 
\end{itemize}

When dealing with oversample clock timed devices, the only possibility is the reference clock synchronization. First of all, the ADCs on the devices are reset and configured to generate samples after the same number of clock cycles. Then, the acquisition is started synchronously with a shared trigger signal. During the reset and sampling procedures, all the devices are locked to the 100~MHz clock signal provided by the PXI chassis. 

\paragraph{Conclusions} National Instruments solutions are, of course, the best synchronization mechanisms that I found during the research, offering extremely high performances and a great software support for easy configuration. All of this comes at certain costs: 
\begin{itemize}
    \item All the modules should be physically connected to the chassis, and no wireless connection is possible (it wouldn't support this kind of synchronization)
    \item The modules should support the National Instruments standard. Even if it is open and can be implemented by manufacturers quite easily, it is targeted at high-cost, high-precision measurement tools.
    \item The PXI system is extremely expensive: the cheapest controller has a cost that is around 1400€ (PXIe-1090), with no module installed. 
\end{itemize} 